% GENERATED FILE. Edit canonical lessons, then run npm run generate:books.
% canonical-source-hash: fnv1a64:08767635778849e6
% canonical-lessons: TE-C486-kolata, TE-C486-parimanam, TE-C486-namuna, TE-C486-viluva, TE-C486-tucu

\chapter{Weights and Measures}
\label{ch:te-weights-and-measures}
\hlchaptermodality{486}

\begin{chapteropening}
\textbf{By the end of this chapter:} \emph{I can say a measurement, size, a sample, value and to weigh.}

Complete the last lesson of chapter 486: I can say a measurement, size, a sample, value and to weigh.
\end{chapteropening}

\section[\texorpdfstring{\te{కొలత} (kolata)}{కొలత (kolata)}]{\te{కొలత} (kolata) — a measurement}

\label{lesson:TE-C486-kolata}



\begin{quote}
Before the new one: say the Telugu for an exchange rate, then the Telugu for to dodge.
\end{quote}

\subsection*{You'll want to know: \te{కొలత}}
\textbf{\te{కొలత}} — \emph{kolata} — \textquotedblleft{}a measurement\textquotedblright{}.

\begin{quote}
\te{దర్జీ} \te{నా} \te{చొక్కా} \te{కొలతలు} \te{తీసుకున్నాడు}.
\end{quote}

\begin{quote}
\emph{darjī nā cokkā kolatalu tīsukunnāḍu.}
\end{quote}

\begin{quote}
The tailor took my measurements for a shirt.
\end{quote}

Plural: \te{కొలతలు}, \emph{kolatalu}.

The noun of \te{కొలుచు}, to measure. \te{కొలతలు} \te{తీసుకో}, to take measurements.

\begin{cousinweb}
\textbf{\te{కొలత}} (\emph{kolata}) is native Dravidian: DEDR 1827 lists Telugu \emph{kolãta}, \textquotedblleft{}measurement, measuring\textquotedblright{}, beside \emph{kolucu}, \textquotedblleft{}to measure\textquotedblright{}.
\end{cousinweb}

\begin{grammarlens}[title={What you've built}]
One more word for getting around.
\end{grammarlens}

\subsection*{Your turn}
\begin{itemize}
\raggedright
  \item \emph{Say it:} \emph{kolata}
  \item \emph{Say it:} \emph{kolata}, once more
  \item \emph{Say it:} \emph{darjī nā cokkā kolatalu tīsukunnāḍu.} — The tailor took my measurements for a shirt.
  \item \emph{Recall:} say the Telugu for an exchange rate, then the Telugu for to dodge, then say \emph{kolata} again
\end{itemize}

\begin{tcolorbox}[breakable,colback=teal!4,colframe=teal!35!black,title={Before you move on}]
What does \te{కొలత} mean? (\textquotedblleft{}A measurement\textquotedblright{}.) Say it once more.
\end{tcolorbox}
\section[\texorpdfstring{\te{పరిమాణం} (parimāṇaṁ)}{పరిమాణం (parimāṇaṁ)}]{\te{పరిమాణం} (parimāṇaṁ) — size; quantity}

\label{lesson:TE-C486-parimanam}



\begin{quote}
Before the new one: say the Telugu for to dodge, then the Telugu for a measurement.
\end{quote}

\subsection*{You'll want to know: \te{పరిమాణం}}
\textbf{\te{పరిమాణం}} — \emph{parimāṇaṁ} — \textquotedblleft{}size; quantity\textquotedblright{}.

\begin{quote}
\te{ఆ} \te{పెట్టె} \te{పరిమాణం} \te{చాలా} \te{పెద్దది}.
\end{quote}

\begin{quote}
\emph{ā peṭṭe parimāṇaṁ cālā peddadi.}
\end{quote}

\begin{quote}
That box is very big in size.
\end{quote}

Size of a thing, or an amount: \te{నీటి} \te{పరిమాణం}, the quantity of water.

\begin{cousinweb}
\textbf{\te{పరిమాణం}} (\emph{parimāṇaṁ}) is borrowed from Sanskrit \emph{parimāṇa}, \textquotedblleft{}measure, size, quantity\textquotedblright{} (Monier-Williams), from \emph{pari-mā}, \textquotedblleft{}to measure out\textquotedblright{}.
\end{cousinweb}

\begin{grammarlens}[title={What you've built}]
One more word for getting around.
\end{grammarlens}

\subsection*{Your turn}
\begin{itemize}
\raggedright
  \item \emph{Say it:} \emph{parimāṇaṁ}
  \item \emph{Say it:} \emph{parimāṇaṁ}, once more
  \item \emph{Say it:} \emph{ā peṭṭe parimāṇaṁ cālā peddadi.} — That box is very big in size.
  \item \emph{Recall:} say the Telugu for to dodge, then the Telugu for a measurement, then say \emph{parimāṇaṁ} again
\end{itemize}

\begin{tcolorbox}[breakable,colback=teal!4,colframe=teal!35!black,title={Before you move on}]
What does \te{పరిమాణం} mean? (\textquotedblleft{}Size; quantity\textquotedblright{}.) Say it once more.
\end{tcolorbox}
\section[\texorpdfstring{\te{నమూనా} (namūnā)}{నమూనా (namūnā)}]{\te{నమూనా} (namūnā) — a sample; a pattern}

\label{lesson:TE-C486-namuna}



\begin{quote}
Before the new one: say the Telugu for a measurement, then the Telugu for size.
\end{quote}

\subsection*{You'll want to know: \te{నమూనా}}
\textbf{\te{నమూనా}} — \emph{namūnā} — \textquotedblleft{}a sample; a pattern\textquotedblright{}.

\begin{quote}
\te{దుకాణదారు} \te{బట్ట} \te{నమూనా} \te{చూపించాడు}.
\end{quote}

\begin{quote}
\emph{dukāṇadāru baṭṭa namūnā cūpiñcāḍu.}
\end{quote}

\begin{quote}
The shopkeeper showed a sample of the cloth.
\end{quote}

A sample of goods, and a pattern or model to copy.

\begin{cousinweb}
\textbf{\te{నమూనా}} (\emph{namūnā}) is a loan from Persian \emph{namūna}, \textquotedblleft{}a sample, a specimen, a model\textquotedblright{}, which reached Telugu through Hindustani.
\end{cousinweb}

\begin{grammarlens}[title={What you've built}]
One more word for getting around.
\end{grammarlens}

\subsection*{Your turn}
\begin{itemize}
\raggedright
  \item \emph{Say it:} \emph{namūnā}
  \item \emph{Say it:} \emph{namūnā}, once more
  \item \emph{Say it:} \emph{dukāṇadāru baṭṭa namūnā cūpiñcāḍu.} — The shopkeeper showed a sample of the cloth.
  \item \emph{Recall:} say the Telugu for a measurement, then the Telugu for size, then say \emph{namūnā} again
\end{itemize}

\begin{tcolorbox}[breakable,colback=teal!4,colframe=teal!35!black,title={Before you move on}]
What does \te{నమూనా} mean? (\textquotedblleft{}A sample; a pattern\textquotedblright{}.) Say it once more.
\end{tcolorbox}
\section[\texorpdfstring{\te{విలువ} (viluva)}{విలువ (viluva)}]{\te{విలువ} (viluva) — value; worth}

\label{lesson:TE-C486-viluva}



\begin{quote}
Before the new one: say the Telugu for size, then the Telugu for a sample.
\end{quote}

\subsection*{You'll want to know: \te{విలువ}}
\textbf{\te{విలువ}} — \emph{viluva} — \textquotedblleft{}value; worth\textquotedblright{}.

\begin{quote}
\te{ఆ} \te{పాత} \te{నాణేలకు} \te{ఇప్పుడు} \te{చాలా} \te{విలువ} \te{ఉంది}.
\end{quote}

\begin{quote}
\emph{ā pāta nāṇēlaku ippuḍu cālā viluva undi.}
\end{quote}

\begin{quote}
Those old coins are worth a lot now.
\end{quote}

Value in money and in esteem: \te{మాటకు} \te{విలువ}, the worth of a promise.

\begin{cousinweb}
\textbf{\te{విలువ}} (\emph{viluva}) is native Dravidian: DEDR 5421 lists Telugu \emph{viluva}, \textquotedblleft{}price, value, cost\textquotedblright{}, beside \emph{vilucu}, \textquotedblleft{}to sell, to buy\textquotedblright{}.
\end{cousinweb}

\begin{grammarlens}[title={What you've built}]
One more word for getting around.
\end{grammarlens}

\subsection*{Your turn}
\begin{itemize}
\raggedright
  \item \emph{Say it:} \emph{viluva}
  \item \emph{Say it:} \emph{viluva}, once more
  \item \emph{Say it:} \emph{ā pāta nāṇēlaku ippuḍu cālā viluva undi.} — Those old coins are worth a lot now.
  \item \emph{Recall:} say the Telugu for size, then the Telugu for a sample, then say \emph{viluva} again
\end{itemize}

\begin{tcolorbox}[breakable,colback=teal!4,colframe=teal!35!black,title={Before you move on}]
What does \te{విలువ} mean? (\textquotedblleft{}Value; worth\textquotedblright{}.) Say it once more.
\end{tcolorbox}
\section[\texorpdfstring{\te{తూచు} (tūcu)}{తూచు (tūcu)}]{\te{తూచు} (tūcu) — to weigh}

\label{lesson:TE-C486-tucu}



\begin{quote}
Before the new one: say the Telugu for a sample, then the Telugu for value.
\end{quote}

\subsection*{You'll want to know: \te{తూచు}}
\textbf{\te{తూచు}} — \emph{tūcu} — \textquotedblleft{}to weigh\textquotedblright{}.

\begin{quote}
\te{కొట్టువాడు} \te{బియ్యం} \te{తూచి} \te{ఇచ్చాడు}.
\end{quote}

\begin{quote}
\emph{koṭṭuvāḍu biyyaṁ tūci iccāḍu.}
\end{quote}

\begin{quote}
The shopkeeper weighed out the rice for us.
\end{quote}

Past (I …): \te{తూచాను}, \emph{tūcānu}.

\te{తూకం}, weight, is its noun. Also to weigh words: \te{మాటలు} \te{తూచి} \te{మాట్లాడు}, to speak carefully.

\begin{cousinweb}
\textbf{\te{తూచు}} (\emph{tūcu}) is native Dravidian: DEDR 3376 lists Telugu \emph{tūcu}, \textquotedblleft{}to weigh, to find the weight of\textquotedblright{}, beside \emph{tūgu}, \textquotedblleft{}to hang, to swing\textquotedblright{}.
\end{cousinweb}

\begin{grammarlens}[title={What you've built}]
One more word for getting around.
\end{grammarlens}

\subsection*{Your turn}
\begin{itemize}
\raggedright
  \item \emph{Say it:} \emph{tūcu}
  \item \emph{Say it:} \emph{tūcu}, once more
  \item \emph{Say it:} \emph{koṭṭuvāḍu biyyaṁ tūci iccāḍu.} — The shopkeeper weighed out the rice for us.
  \item \emph{Recall:} say the Telugu for a sample, then the Telugu for value, then say \emph{tūcu} again
\end{itemize}

\begin{tcolorbox}[breakable,colback=teal!4,colframe=teal!35!black,title={Before you move on}]
What does \te{తూచు} mean? (\textquotedblleft{}To weigh\textquotedblright{}.) Say it once more.
\end{tcolorbox}
